\section{Background: why this is an Information Systems problem}
\label{sec:background}

\subsection{The rule, the specification and the code}

In the IS tradition an information system is understood as a system of work in which human
participants and technological components jointly perform activities to process information and
support decision-making, and its value emerges from integration into organisational practice rather
than from technological functionality alone. The alignment between business strategy and information
technology has been framed as a continuing process of managerial fit rather than a one-off design
choice \citep{henderson1993strategic}, and the success of an information system has been modelled as a
multi-dimensional construct in which the quality of the information and the quality of the system are
distinct from their use and their net organisational benefit \citep{delone2003model}. These framings
matter here for a direct reason: they insist that a system can be technically sound and
organisationally wrong at the same time, which is precisely the failure mode this paper measures.

BDD is an attempt to make that correspondence inspectable. A requirement is expressed as scenarios in
Gherkin, in domain vocabulary, and those scenarios are bound to executable steps. Three artefacts are
then in play, and each can move independently:

\begin{itemize}
\item the \textbf{organisational rule} (what the business intends, possibly recorded only in a policy
      document or a verbal agreement);
\item the \textbf{executable specification} (the Gherkin scenarios, which are supposed to state the
      rule in checkable form);
\item the \textbf{implementation} (the code that actually decides outcomes).
\end{itemize}

BDD's premise is that the second artefact binds the first and the third: change the code without
changing the scenarios and the suite fails. The premise fails quietly when the second artefact loses
its discriminating power. A scenario whose step definition asserts nothing, or asserts that a value
equals itself, or compares the actual output with a value recomputed using the same production code,
will not fail when the implementation changes --- because the oracle moved with the implementation.
The specification is then \emph{descriptively} present and \emph{normatively} absent.

\subsection{Why coverage cannot detect this}

The IS literature on system evaluation is familiar with the difference between \emph{utilisation} and
\emph{effectiveness}. Test coverage has the same character: it records that code was exercised, not
that exercising it could have produced a failure. The empirical record is unambiguous. In an
industrial Ruby project, mutation scores averaged 0.78 while statement coverage averaged 97.2\,\%
\citep{li2015mutationpracticeruby}. Across 45 Maven projects, the correlation between coverage and
mutation score was 0.75 when computed na\"ively but fell to 0.21 once restricted to covered code
\citep{Jain2023MindTheGap} --- that is, almost all of the apparent association was an artefact of
\emph{which} code was covered, and among code that is covered, coverage says very little about
discrimination. Mutation score, by contrast, tracks real fault detection independently of coverage
\citep{Just2014AreMutants}, although it remains a proxy: a substantial minority of real faults have no
mutant representative \citep{Just2014AreMutants,titcheuchekam2017faultrevelation}.

For an organisation the practical question is therefore not ``what is our coverage?'' but ``if
someone changed this business rule, would anything fail?'' That is the question mutation analysis
answers, and it is the question this paper operationalises for the specification layer specifically.

\subsection{Decay is the normal case, not the exception}

That test artefacts degrade is well documented. Disabled tests are common --- 3{,}111 disabling changes
across 122{,}000 commits in 15 Java systems, with 41\,\% of disabled tests never re-enabled and a
median 39 days to re-enable \citep{Kim2021Disabled}. Test cases have a half-life of five to twelve
months in a large industrial setting \citep{Feldt2014GrowOld}. A substantial share of ``green'' tests
check less than their authors believe: 294 rotten green tests were found in 19{,}905 test cases
\citep{Delplanque2019RottenGreen}, and inadequate assertions have been catalogued as a recurring
defect pattern \citep{zhi2013inadequateassertions,zhang2015assertions}. Test smells are pervasive ---
86\,\% of JUnit tests in one study exhibited at least one \citep{Bavota2015AreTestSmells} --- and their
removal is mostly incidental rather than deliberate \citep{Kim2021SecretLife}.

If decay is the normal case for tests in general, there is no reason to expect the specification layer
of a BDD suite to be exempt. It is, however, a layer with a distinctive failure mode: because the
scenario text remains readable and business-facing even when the assertions behind it have been
hollowed out, the artefact continues to \emph{look} like a specification long after it has stopped
functioning as one. That asymmetry --- legible appearance, degraded function --- is what makes an
instrument necessary, and it is why a practitioner cannot simply read the feature files and know.

\subsection{The Design Science framing}

We therefore treat the problem as a design problem, in the sense of Design Science Research (DSR).
Design science builds and evaluates artefacts intended to address identified organisational problems,
with relevance and rigour as twin requirements \citep{hevner2004design}, and it is distinguished from
natural-science research by the kind of knowledge it produces: the one builds artefacts, the other
seeks explanations, and the two have different truth criteria \citep{march1995design}. DSR is
accordingly expected to state the problem class, the artefact's design rationale, its evaluation in a
relevant context, and the implications for practice \citep{peffers2007dsrm,gregor2013positioning}. The
response here is to build and evaluate an artefact that makes the divergence detectable.

Framing the contribution this way is not a rhetorical choice. It determines what counts as success. A
DSR artefact is evaluated on whether it solves the problem class in a relevant context, and a
negative finding about the context --- that the population on which the artefact would be validated at
scale barely exists --- is a finding about the artefact's deployment conditions, not a failure of the
artefact. Section~\ref{sec:discussion} draws out that distinction.
